\documentclass[12pt,a4paper]{article}
\usepackage[utf8]{inputenc}
\usepackage{amsmath,amssymb,amsthm}
\usepackage{algorithm}
\usepackage{algpseudocode}
\usepackage{geometry}
\usepackage{hyperref}
\usepackage{booktabs}

\geometry{margin=2.5cm}

\newtheorem{theorem}{Theorem}
\newtheorem{lemma}[theorem]{Lemma}
\newtheorem{corollary}[theorem]{Corollary}
\newtheorem{definition}{Definition}

\newcommand{\Fmn}{F(m,n)}
\newcommand{\dn}{d_n}
\newcommand{\MOD}{1\,234\,567\,891}

\title{Problem 452: Long Products}
\author{}
\date{}

\begin{document}
\maketitle

\section{Problem Statement}

Define $\Fmn$ as the number of $n$-tuples $(a_1, a_2, \ldots, a_n)$ of positive
integers for which the product of the elements does not exceed~$m$:
\[
  \Fmn = \bigl|\bigl\{(a_1, \ldots, a_n) \in \mathbb{Z}_{>0}^n : a_1 a_2 \cdots a_n \le m\bigr\}\bigr|.
\]

Given that $F(10, 10) = 571$ and $F(10^6, 10^6) \bmod \MOD = 252\,903\,833$,
find $F(10^9, 10^9) \bmod \MOD$.

\section{Mathematical Foundation}

\begin{theorem}[Reduction to Piltz divisor sum]\label{thm:piltz}
Grouping tuples by product value,
\[
  \Fmn = \sum_{k=1}^{m} \dn(k),
\]
where $\dn(k)$ counts the ordered factorizations of~$k$ into exactly~$n$ positive integer factors.
\end{theorem}

\begin{proof}
Each $n$-tuple $(a_1, \ldots, a_n)$ with product $k$ is an ordered factorization of $k$ into $n$ parts. Summing over all $k \le m$ counts every valid tuple exactly once.
\end{proof}

\begin{lemma}[Multiplicativity of $\dn$]\label{lem:mult}
The function $\dn$ is multiplicative: for $\gcd(a,b) = 1$,
$\dn(ab) = \dn(a)\,\dn(b)$. At a prime power $p^a$:
\[
  \dn(p^a) = \binom{n + a - 1}{a}.
\]
\end{lemma}

\begin{proof}
Multiplicativity follows from CRT applied to ordered factorizations. For $p^a$, distributing $a$ copies of prime~$p$ among $n$ positions is a stars-and-bars count: $\binom{n+a-1}{a}$.
\end{proof}

\begin{lemma}[Exponent bound]\label{lem:exp}
Since $m = 10^9$ and $2^{30} > 10^9$, every prime power $p^a \le m$ satisfies $a \le 29$. Hence $\binom{n+a-1}{a}$ is a polynomial of degree at most~29 in~$n$, readily evaluable modulo the prime $\MOD$.
\end{lemma}

\begin{proof}
If $p^a \le 10^9$ then $a \le \lfloor \log_2(10^9) \rfloor = 29$. The binomial coefficient $\binom{n+a-1}{a} = \prod_{j=0}^{a-1} (n+j) / a!$ is a polynomial of degree~$a$ in~$n$.
\end{proof}

\begin{theorem}[Bucket decomposition]\label{thm:bucket}
The set $\{\lfloor m/k \rfloor : k = 1, \ldots, m\}$ has at most $2\lfloor\sqrt{m}\rfloor$ distinct values. Storing partial sums $S[v]$ for each such~$v$ suffices to compute $\Fmn$ with $O(\sqrt{m})$ storage.
\end{theorem}

\begin{proof}
For $k \le \sqrt{m}$, the value $\lfloor m/k \rfloor$ ranges over at most $\sqrt{m}$ values. For $k > \sqrt{m}$, $\lfloor m/k \rfloor < \sqrt{m}$, contributing at most $\sqrt{m}$ additional values. The union has at most $2\sqrt{m}$ elements.
\end{proof}

\section{Algorithm}

\begin{algorithm}[H]
\caption{Compute $\Fmn \bmod P$ via multiplicative bucket sieve}
\begin{algorithmic}[1]
\State $B \gets \lfloor\sqrt{m}\rfloor$
\State Initialize $\texttt{lo}[v] \gets 1$ for $v = 1, \ldots, B$ and $\texttt{hi}[k] \gets 1$ for $k = 1, \ldots, B$
\State Precompute $C[a] \gets \binom{n+a-1}{a} \bmod P$ for $a = 0, \ldots, 29$
\For{each prime $p \le B$} \Comment{Small primes}
    \For{each bucket value $v$ in decreasing order}
        \For{$a = 1, 2, \ldots, \lfloor\log_p v\rfloor$}
            \State $S[v] \mathrel{+}= C[a] \cdot S[\lfloor v/p^a \rfloor] \pmod{P}$
        \EndFor
    \EndFor
\EndFor
\For{each prime $p \in (B, m]$ via segmented sieve} \Comment{Large primes}
    \For{$k = 1$ to $\min(B, \lfloor m/p \rfloor)$}
        \State $\texttt{hi}[k] \mathrel{+}= C[1] \cdot \texttt{lo}[\lfloor \lfloor m/k \rfloor / p \rfloor] \pmod{P}$
    \EndFor
\EndFor
\State \Return $\texttt{hi}[1]$
\end{algorithmic}
\end{algorithm}

\section{Complexity Analysis}

\begin{itemize}
    \item \textbf{Time:} $O(m)$, dominated by the large-prime phase. The small-prime phase costs $O(\sqrt{m} \cdot \sqrt{m} / \log \sqrt{m}) = O(m / \log m)$. Each large prime~$p$ updates at most $\lfloor m/p \rfloor \le \sqrt{m}$ buckets; summing gives $O(\sum_{p > \sqrt{m}} m/p) = O(m)$.
    \item \textbf{Space:} $O(\sqrt{m})$ for the two bucket arrays and the segmented sieve buffer.
\end{itemize}

\section{Answer}

\[
\boxed{345558983}
\]

\end{document}
